\PDEIndexMark{idx:banach-inverse-function-theorem-introduction}关于方程可解性的关键结果是逆函数定理, 它说明若方程的线性化局部唯一可解, 则该方程也局部唯一可解.

\begin{Theorem}{逆函数定理}
\label{thm:banach-inverse-function}
设 $X$ 和 $Y$ 为 Banach 空间, $U$ 为 $X$ 中原点的开邻域. 设 $F:U\to Y$ 连续可微, 并假设 $F^{\prime}(0)\in\mathcal{L}(X,Y)$ 既单又满. 则存在 $Y$ 中 $F(0)$ 的开邻域 $V$ 及连续可微映射 $G:V\to X$, 使 $F(G(y))=y$. 此外, $G(y)$ 是方程 $F(x)=y$ 的唯一足够小的解 $x$.
\end{Theorem}

\begingroup
\renewcommand{\endProof}{\par\nobreak\hfill$\blacksquare$\par\vspace{0.1cm}}
\begin{Proof}
考虑由
\begin{equation}\label{eq:inverse-function-contraction-map-definition}
T_y(x)=F^{\prime}(0)^{-1}(y-F(0))-F^{\prime}(0)^{-1}(F(x)-F(0)-F^{\prime}(0)x).
\end{equation}
定义的映射 $T_y$.令 $M=\|F^{\prime}(0)^{-1}\|$. 断言: 若 $\epsilon$ 选得足够小, 且 $\|y-F(0)\|<\epsilon/(2M)$, 则 $T_y$ 为将 $X$ 中闭 $\epsilon$ 球映入自身的压缩映射. 为此, 注意
\begin{equation}\label{eq:inverse-function-contraction-ball-invariance-bound}
\|T_y(x)\|\leq M\|y-F(0)\|+M\|F(x)-F(0)-F^{\prime}(0)x\|.
\end{equation}
右端第一项小于 $\epsilon/2$, 而由导数的定义, 第二项为 $o(\epsilon)$. 因此, 若 $\epsilon$ 足够小, 则有 $\|T_y(x)\|<\epsilon$. 此外, 注意
\begin{equation}\label{eq:inverse-function-contraction-map-derivative}
T_y^{\prime}(x)=-F^{\prime}(0)^{-1}(F^{\prime}(x)-F^{\prime}(0))
\end{equation}
在 $\epsilon$ 足够小时的范数小于例如 $1/2$. 压缩性质现在由等式
\begin{equation}\label{eq:inverse-function-contraction-integrated-derivative}
T_y(x)-T_y(z)=\int_0^1T_y^{\prime}(z+\theta(x-z))(x-z)\,\mathrm{d}\theta.
\end{equation}
得到.由 Banach 压缩原理 (Theorem~\ref{thm:banach-contraction-fixed-point}), $T_y$ 在半径为 $\epsilon$ 的球中有唯一不动点. 将此不动点记为 $G(y)$. 容易验证, 方程 $T_y(x)=x$ 等价于 $F(x)=y$.

还需证明 $G$ 实际上连续可微. 先证明连续性. 有 $G(y)=T_y(G(y))$, 因而
\begin{equation}\label{eq:inverse-function-lipschitz-continuity-estimate}
\begin{aligned}
\|G(y+h)-G(y)\|&\leq\|T_{y+h}(G(y+h))-T_y(G(y+h))\|\\
&\quad+\|T_y(G(y+h))-T_y(G(y))\|\\
&\leq M\|h\|+\frac{1}{2}\|G(y+h)-G(y)\|,
\end{aligned}
\end{equation}
其中最后一个估计来自 $T_y$ 的压缩性质. 由 \eqref{eq:inverse-function-lipschitz-continuity-estimate} 立即可知 $G$ 为 Lipschitz 连续.

现在令 $k=G(y+h)-G(y)$, 则由 \eqref{eq:inverse-function-lipschitz-continuity-estimate} 立即得到 $\|k\|\leq2M\|h\|$. 有
\begin{equation}\label{eq:inverse-function-differentiability-remainder-expansion}
0=F(G(y+h))-F(G(y))-h=F^{\prime}(G(y))k+R(k)-h,
\end{equation}
其中当 $\|k\|\to0$ 时, $\|R(k)\|/\|k\|\to0$. 从而
\begin{equation}\label{eq:inverse-function-increment-linearization}
k=F^{\prime}(G(y))^{-1}h-F^{\prime}(G(y))^{-1}R(k),
\end{equation}
又由于 $\|k\|\leq2M\|h\|$, 当 $\|h\|\to0$ 时有 $\|R(k)\|/\|h\|\to0$. 这证明 $G$ 可微, 且导数为 $F^{\prime}(G(y))^{-1}$.
\end{Proof}
\endgroup

\begin{Remark}
\label{rem:inverse-function-higher-differentiability}
若 $F$ 为 $k$ 次连续可微, 则 $G$ 也如此 (Problem~\ref{ex:inverse-function-higher-differentiability}).
\end{Remark}
